% 第 44 章　薛定谔方程：量子态如何随时间变化
\chapter{薛定谔方程：量子态如何随时间变化}

第~42~章我们看到经典物理在黑体辐射、光电效应和原子光谱面前连输三阵；第~43~章我们学会了新的记账工具：量子态用波函数表示，振幅的模平方给出概率，复数负责记录相位。但那一章的账簿是静态的——它只告诉你“此刻测量会得到什么”。一个完整理论还必须回答全书反复追问的第二个问题：\emph{什么规律决定状态如何变化？}牛顿力学用 \(F=ma\) 回答它，量子力学用一条以“能量”为发动机的方程回答它。这条方程就是本章的主角。读完本章，你将明白为什么同一种材料磨成不同大小的颗粒会发出不同颜色的光，为什么微观世界存在“永不变调”的音符，以及你的 U 盘为什么终有一天会忘掉里面的数据。

\step{反常识开场}

先看一种你或许已经买回家的东西。高端液晶电视的广告里常出现“量子点技术”：屏幕里掺着亿万个纳米级的小颗粒，成分是同一种半导体材料（比如硒化镉）。奇妙的事情是：把这些颗粒磨成直径两三纳米，它们发蓝光；磨成五六纳米，发绿光；再大一点，发红光。\emph{化学成分表一模一样，唯一的变量是颗粒的大小}。我们从小被教导：颜色由材料决定——黄金是黄的，铜是红的，硫酸铜溶液是蓝的。可在纳米尺度上，材料把决定权交了出去：颜色由“盒子有多大”定价。颜色是光的频率，频率在量子世界里对应能量——也就是说，一颗电子的能量，居然能被它“住的房子”的尺寸强行规定。这没有任何经典对照物：你把一个弹珠关进大盒子还是小盒子，它的能量完全由你扔得多用力决定，盒子尺寸无权置喙。

再看一个相反方向的现象：不是“变”，而是“惊人的不变”。氢原子受激后会发出一种特有的红光，波长在 \(656.3\) 纳米附近。你在实验室点亮一盏氢灯，它是这个波长；天文望远镜对准一百多亿光年外的星系，那里氢原子发出的光，红移修正之后，小数点后的每一位都完全相同。今天的物理学家把这种稳定性榨到了极限：锶原子钟利用原子内部一次跃迁的频率计时，不确定度小于 \(10^{-18}\)——如果让它从宇宙大爆炸走到今天再走一遍，累计误差不到一秒。宏观世界里不存在永不变调的东西：琴弦会松，发条会软，晶体振荡器会老化。微观世界却有一类状态，仿佛被时间冻结，永远走音不了。

一个现象说“尺寸能规定能量”，一个现象说“有些状态能被时间冻结”。要同时解释它们，就必须回答本章的问题：量子态到底怎样随时间变化？

\step{先猜一猜}

老规矩，先押注。

第一题。把一颗电子关进一个无法逃脱的盒子里。现在把盒子的宽度压缩成原来的一半，这颗电子的能量会（a）不变；（b）变成两倍；（c）变成四倍；（d）变小？

第二题。你朝一个小土坡滚一个球，球的速度不够，滚到半山腰必然退回——这是经典世界的铁律：能量不够，就是过不去。现在换成一颗电子，面对一堵“能量不够”的势能高墙，它有没有可能出现在墙的另一边？（a）绝无可能，铁律就是铁律；（b）有一点点可能；（c）一定能穿过去。

第三题。量子力学里有一类特殊状态叫“定态”，它的概率分布不随时间变化。一颗处于定态的电子，是（a）静止在某个位置上；（b）仍然“不安分”，动能不为零；（c）概率不动就等于一切都不动？

第四题。把两个“各自永远不变”的定态按一定比例叠加起来，得到的新状态（a）当然也永远不变——不变加不变还是不变；（b）会随时间来回振荡。请认真想想再写，这道题是专门给直觉准备的陷阱。

写下你的四个答案，本章结束时逐一对账。

\step{动手实验}

\begin{experiment}[一根橡皮筋的音阶]
\textbf{材料}：一根粗橡皮筋（或吉他、尤克里里的任意一根弦）、一张桌子、一部手机（有慢动作摄像更好）。

\textbf{步骤一}：把橡皮筋拉长，两端分别固定在桌子两端（用胶带、书本压住均可），绷到拨一下能发出清晰音调的程度。拨动它的正中间，听音调，看它振动——整根皮筋一起上下，呈一个饱满的拱形。

\textbf{步骤二}：用一根手指\emph{轻轻}点在皮筋的正中央（不要压死，只是贴住），再拨动其中一半。皮筋发出高了一个八度的音。仔细观察或用手机慢动作拍摄：皮筋此刻以两段反向的拱在振动，中点始终不动。再把手指轻点在三分之一处拨动，音调又变了，皮筋呈现三段起伏。

\textbf{步骤三}：把一端的固定点挪近，让振动长度缩短一半，直接拨动整根：音调大约翻倍。

\textbf{步骤四}：试试“乱来”——比如在 \(37\%\) 的位置使劲压出一个任意的波形松手。松手瞬间波形立刻散开、晃动、很快重组为混乱的抖动，得不到任何干净的音调。

\textbf{记录与思考}：为什么只有“整数段”的波形能稳定驻留，任意波形却活不过一瞬间？弦长和音调是什么关系？请把这两个观察记牢，并且做个思想替换：把“弦长”想成电子的“盒子”，把“音调”想成电子的“能量”——你手里这根橡皮筋，就是本章主角方程的宏观预告片。
\end{experiment}

\step{旧模型遇到麻烦}

要找一个新的演化规律，先得说清楚旧规律为什么不能用。我们的候选人一个个上庭。

\textbf{候选人一：牛顿第二定律。}它的“现在”是两个数：位置和速度；给定它们，力就接管未来。但第~43~章已经判了这幅图景死刑：电子没有偷偷摸摸的确定轨迹，双缝条纹是铁证。量子态的全部内容是整张波函数地图——所以新定律的输入不能是两个数，必须是整张地图本身。牛顿方程连“输入格式”都对不上，出局。

\textbf{候选人二：经典波动方程。}既然电子表现出波动性，能不能直接借用绳波、水波的方程？这类方程的骨架是“位移的变化的变化（加速度），正比于位移的弯曲程度”——它对时间是\emph{二阶}的，而且是\emph{实数}的。两个麻烦立刻出现。其一，波函数不是可以围观的位移，它是概率幅；其二，更致命的是定态难题：本章开场那些“概率分布被时间冻结”的状态，要求振幅的大小不变、内部却还在演化（否则无法解释两个定态叠加后会随时间振荡，这正是先猜一猜的第四题）。一个实数的波，幅度不变就是什么都不变，做不到“外形冻结、内部旋转”。要让“变化”只发生在相位上，必须请出复数——第~43~章引入复数时的理由，在这里从“方便”升级成了“必需”。

\textbf{候选人三：只用能量守恒。}“知道系统总能量”能不能推出状态怎么变？不能。同一个总能量可以对应无数种完全不同的状态：盒中电子可以安静地待在最低的驻波上，也可以是好几个驻波的叠加。总能量是一个数字，而状态的演化需要知道能量在空间里如何\emph{分布}——哪一份是“运动的”（动能），哪一份是“位置的”（势能），这个分工随地点不同而不同。一个总数管不了这么细。

旧候选人全部出局后，新定律必须满足的硬约束反而清晰了，一共三条。第一，概率总和必须永远是一——演化不能让概率凭空变多或变少。第二，叠加原理必须成立——两个可演化出的状态，它们的叠加也要能演化，且演化结果就是各自演化结果的叠加（否则第~43~章的干涉账就记不下去）。第三，宏观极限必须回到牛顿——用大质量的、位置相当明确的波包去算，它中心的移动必须服从 \(F=ma\)，否则整个经典大厦就塌了。薛定谔在一九二六年初写下方程时，手里握的正是这几条约束、德布罗意的波长关系，以及一个关键启发：经典力学里的哈密顿量——那本把动能和势能分列两栏的能量总账——早就偷偷兼任着“时间演化发电机”的职务。他要做的，只是把这台发电机翻译成量子语言。

\step{发明新概念}

\begin{mainline}
\textbf{第一个新概念：哈密顿量——能量账本，兼任时间发电机。}经典力学里，一个系统的能量分两类记：动能 \(\tfrac12 mv^{2}\) 管“它动得多欢”，势能 \(V\) 管“它待的位置值多少钱”。把这两栏加起来，写成 \(H\)，叫哈密顿量。它在经典力学里就有一项隐藏职务：知道 \(H\)，就能推出系统下一步怎么动。量子力学把这项职务扶正：\emph{哈密顿量是演化的唯一发动机}。它怎么“发动”？对着波函数干活：动能那一栏关心波函数在空间上\emph{弯得有多急}（弯曲对应波长，波长对应动量，动量平方对应动能），势能那一栏关心波函数\emph{站在什么位置}。两栏合起来，决定了波函数下一瞬间往哪个方向变。注意它已经从“一个数”升级成了“一种作用在状态上的操作”——这类操作叫算符，第~46~章会系统展开，本章只需要记住：哈密顿量读出的是能量，生成的是变化。

\textbf{第二个新概念：薛定谔方程（概念版）。}整条方程用一句话说：\emph{波函数的变化率，正比于哈密顿量作用在波函数上的结果}，比例系数里带一个 \(i\) 和一个很小的常数 \(\hbar\)（读作“h 拔”，它等于普朗克常数 \(h\) 除以 \(2\pi\)）。每一项都有不得不存在的理由。哈密顿量在场，因为演化必须由能量驱动——这是从经典力学继承的发电机。\(i\) 在场，因为它让变化变成\emph{旋转}而不是\emph{增减}：乘以 \(i\) 等于把复数箭头转过九十度，反复转就是绕圈——箭头长度（概率）分毫不变，方向（相位）时刻在变。这正是“概率总和恒为一”的数学保障。\(\hbar\) 在场，负责把能量翻译成“相位转速”，并保证宏观世界里量子效应小到看不见。它像什么？像钟表里的发条：能量越大，指针转得越欢。它不像什么？钟表指针是空间里一根看得见的杆，而这里的“转动”发生在每个地点的复数小箭头上，不占据任何空间，也无法被直接围观——你能看到的永远只有它平方之后的概率。

\textbf{第三个新概念：定态——能量确定的状态，被时间冻结的概率地形。}如果波函数恰好是“能量取单一确定值”的形状，薛定谔方程的演化就退化成一件事：整张地图的相位以均匀的速度整体旋转，角速度恰为 \(\omega=E/\hbar\)。整体旋转不影响任何模平方——于是概率分布、一切可测的统计性质全部冻结。这种能量确定的态叫\emph{定态}。请立刻建立防线：定态\emph{不是}“粒子静止不动”。盒中电子的最低定态能量不为零，动能时刻存在，它更像一根不停振动的琴弦：弦上每一点的\emph{振幅包络}不动，但弦本身正在剧烈地抖动。冻结的是地形，不是里面的东西。

\textbf{第四个新概念：能级——边界条件逼出来的离散。}把电子关进宽度为 \(L\) 的盒子，波函数在盒壁处必须归零（墙外绝不可能找到它）。能在盒子里稳定驻留的，只有恰好把整数个半波塞进盒宽的波形——和你的橡皮筋一模一样：一段、两段、三段，没有一段半。波长被边界卡死，动量就被卡死，能量随之被卡死。于是盒中电子的能量只能取一串离散值：\emph{能级}。离散不是谁的规定，是“墙”逼出来的。它与琴弦音阶的相似到此为止，不同也要说清：琴弦是真实的位移，波函数是概率幅；琴弦的音阶大体均匀分布，盒中电子的能级却按平方散开，越高的能级彼此离得越远。

\textbf{第五个新概念：波包——注定要发胖的自由粒子。}没有墙的自由粒子，能量不再离散。一个局域的、大体知道“它在这一带”的状态，必须是许多种波长的叠加——这叫波包。薛定谔方程对每种波长的成分给出不同的相位转速（能量不同，转速不同），于是各成分之间的相对相位随时间错开，原先靠精密配合堆出来的“鼓包”逐渐散架：\emph{波包必然随时间变胖}。这不是波被什么东西撞散了，是纯粹的相位赛跑——没有任何耗散，演化完全可逆。

\textbf{第六个新概念：隧穿——经典禁区里不立即归零的概率。}面对一堵势能高墙，且粒子能量低于墙高：经典判决是“墙内粒子数为零”。薛定谔方程的判决不同：波函数在墙根处不骤降为零，而是在墙内\emph{指数式地衰减}——像声音穿墙，每过一堵厚度就弱一截。关键在“指数”而不是“零”：只要墙足够薄，衰减还没见底，墙的另一侧就还残留着非零的振幅，粒子就有非零的概率出现在“经典禁区”之外。先猜一猜第二题的答案是（b）。请警惕这个名字的误导：没有任何东西在墙里“挖”了一条路，粒子也没有在墙内走某条隐秘的小径；真实的陈述只有一句——墙外波函数不为零，测量就有概率在那里找到整颗的粒子。
\end{mainline}

\step{一句话模型}

\begin{oneline}
量子态是一张带相位的概率地图；薛定谔方程是它的流动法则——系统的能量充当相位的发条，能量越高的成分相位转得越快，概率地形由能量在空间上的分布塑造；能量单一确定的成分，相位只能整体匀速空转，于是概率分布被时间冻结——这就是定态，冻结的是地形，不是里面的运动。
\end{oneline}

\step{数学翻译器}

\begin{mainline}
\textbf{相位的发条。}第~42~章给出过 \(E=hf\)（光的能量与频率），第~43~章给出过 \(p=h/\lambda\)（德布罗意波长）。把两条拼起来，定态相位旋转的角速度为
\[
  \omega=\frac{E}{\hbar},\qquad \hbar=\frac{h}{2\pi}.
\]
它描述什么：能量为 \(E\) 的定态，其相位箭头每秒转过多少弧度。为什么这样写：这是 \(E=hf\) 把“每秒转几圈”换算成“每秒转几弧度”的同一条账。立刻做例行检查。令 \(E=0\)：\(\omega=0\)，相位纹丝不动，状态彻底冻结——能量为零的态连相位都不转，自洽。能量加倍：转速加倍——所以两个定态叠加时，能量高的成分“跑得快”，两者相对相位随时间错开，这正是第四道猜题的机关：两个各自冻结的定态，叠加后概率分布会以频率
\[
  f=\frac{E_{2}-E_{1}}{h}
\]
来回振荡，像两个音高略有差别的音叉叠出的“拍”。不变加不变，居然变出“动”——你的直觉若选了（a），错因是把“每个成分的模平方不变”误当成“叠加后的模平方不变”，而叠加后的概率里含有两个振幅的\emph{交叉项}，交叉项里活着两者的相位差。

\textbf{盒中粒子：把整数个半波塞进盒子。}盒宽 \(L\)，波函数在两端为零，允许的波长满足
\[
  L=\frac{n\lambda}{2},\qquad n=1,2,3,\dots
\]
代入 \(p=h/\lambda\) 得动量 \(p_{n}=nh/(2L)\)，再代入动能 \(E=p^{2}/(2m)\)：
\[
  E_{n}=\frac{n^{2}h^{2}}{8mL^{2}}.
\]
它描述什么：宽度 \(L\) 的盒子里，质量 \(m\) 的粒子只允许拥有的能量清单。何时成立：壁不可穿透、粒子非相对论性。逐项检查。令 \(n=1\)：最低能量 \(E_{1}=h^{2}/(8mL^{2})\)，\emph{不为零}——零点能。被关起来的粒子无法彻底安静：波函数若处处为零，等于“盒里根本没有粒子”，这个“解”因描述空盒子而被开除。这是本章第一个反直觉结论的严格版：囚禁本身就要花能量，关得越紧（\(L\) 越小），这份“不安分”的能量越贵，而且是平方地贵——盒子宽度减半，每个能级变成四倍。第一题选（c）。令 \(L\) 趋于极大：相邻能级的间隔 \(\Delta E\approx(2n+1)E_{1}\) 趋于零，能量清单密成连续谱——盒子的量子性是墙的产物，墙推到无穷远，经典世界回来了。令 \(m\) 很大：能级间隔反比于质量，宏观物体的能级密到任何仪器都分不出——第三条硬约束（宏观回到牛顿）兑现。

\textbf{真实数据估算：量子点的颜色是谁定的价。}取一颗直径 \(3\) 纳米的硒化镉颗粒，近似当成 \(L=3\times10^{-9}\) 米的盒子。半导体中的电子受晶格的集体影响，表现得像一个“有效质量”只有约 \(0.13\,m_{e}\) 的粒子（\(m_{e}=9.11\times10^{-31}\) 千克）。代入：
\[
  E_{1}=\frac{(6.63\times10^{-34})^{2}}{8\times(0.13\times9.11\times10^{-31})\times(3\times10^{-9})^{2}}
  \approx5\times10^{-20}\ \text{J}\approx0.3\ \text{eV}.
\]
第一激发态与基态之差 \(\Delta E=3E_{1}\approx1\ \text{eV}\) 量级；电子和材料里留下的空穴各贡献一份，再叠上硒化镉本身约 \(1.7\ \text{eV}\) 的固有能隙，发光的总能量从小颗粒的蓝绿端滑到大颗粒的红端——与量子点电视的色谱完全一致。决定性的是那条平方反比：\(E_{n}\propto1/L^{2}\)，颗粒直径加倍，全部量子化的能量缩小到四分之一。颜色确实是被“盒子大小”明码标价的。

\textbf{自由波包的发胖。}波包内波长为 \(\lambda\) 的成分，能量 \(E=p^{2}/(2m)=h^{2}/(2m\lambda^{2})\)。不同 \(\lambda\) 的成分相位转速不同，彼此错开，鼓包散开。检查两个极端。只含单一波长：所有点转速一致，波形不变胖——但代价是它均匀铺满全空间，根本不是一个“局域的粒子”。质量趋于极大：转速差 \(\Delta\omega=\Delta E/\hbar\) 被质量压小，波包几乎不散——棒球飞完全场，它的波包展宽完全可以忽略，这就是为什么你从不觉得棒球在“变胖”。

\textbf{隧穿的半定量。}能量 \(E\) 低于势垒 \(V_{0}\) 时，波函数在墙内按指数衰减，每深入一小段就弱一截；透射概率大体随“墙厚”指数下降。做三个例行检查。墙厚趋于零：没有墙，当然全透射——透射概率趋于一，合理。墙厚趋于无穷：透射概率趋于零——经典判决“过不去”被完整收回为极限情形。粒子能量逼近垒高 \(V_{0}\)：衰减变慢，透射显著上升——所以“差一点点就够高”的电子最容易漏过去。指数依赖是工程上的黄金杠杆，后面讲闪存时兑现。
\end{mainline}

\begin{center}
\begin{tikzpicture}[scale=1.0]
  % 无限深势阱
  \draw[very thick] (0,3.4) -- (0,0) -- (4.8,0) -- (4.8,3.4);
  \node[left] at (0,3.2) {墙};
  \node[right] at (4.8,3.2) {墙};
  \node[below] at (2.4,0) {盒宽 \(L\)};
  % n=1
  \draw[thick,blue!70!black] plot[domain=0:4.8,samples=60,smooth] (\x,{0.7+0.55*sin(180*\x/4.8)});
  \node[right,blue!70!black] at (4.9,1.25) {\(n=1\)};
  % n=2
  \draw[thick,red!70!black] plot[domain=0:4.8,samples=80,smooth] (\x,{1.9+0.55*sin(360*\x/4.8)});
  \node[right,red!70!black] at (4.9,2.45) {\(n=2\)};
  % n=3
  \draw[thick,green!50!black] plot[domain=0:4.8,samples=100,smooth] (\x,{3.0+0.4*sin(540*\x/4.8)});
  \node[right,green!50!black] at (4.9,3.0) {\(n=3\)};
  \node at (2.4,-0.9) {只有整数个半波能塞进盒子：波长被墙卡死，能量随之离散};
\end{tikzpicture}
\end{center}

\begin{center}
\begin{tikzpicture}[scale=1.0]
  % 波包扩散
  \draw[thick,blue!70!black] plot[domain=-2.4:2.4,samples=80,smooth] (\x,{1.8*exp(-3*\x*\x)});
  \node[above,blue!70!black] at (0,1.8) {此刻：瘦而高};
  \draw[thick,red!70!black,dashed] plot[domain=-4.5:4.5,samples=100,smooth] (\x,{0.55*exp(-0.3*\x*\x)});
  \node[red!70!black] at (3.2,0.85) {稍后：胖而矮};
  \draw[-{Stealth[length=2.5mm]},thick,gray] (1.5,0.9) -- (2.6,0.6);
  \draw[-{Stealth[length=2.5mm]},thick,gray] (-1.5,0.9) -- (-2.6,0.6);
  \node at (0,-0.8) {自由波包：各波长成分相位转速不同，鼓包必然随时间散开};
\end{tikzpicture}
\end{center}

\begin{center}
\begin{tikzpicture}[scale=1.0]
  % 势垒
  \fill[gray!25] (3,0) rectangle (4.2,2.0);
  \draw[thick] (0,0) -- (8.4,0);
  \draw[thick] (3,0) -- (3,2.0) -- (4.2,2.0) -- (4.2,0);
  \node at (3.6,2.25) {势垒 \(V_{0}>E\)};
  % 入射波
  \draw[thick,blue!70!black] plot[domain=0:3,samples=80,smooth] (\x,{0.9+0.6*sin(480*\x)});
  \node[blue!70!black] at (1.5,1.9) {入射波};
  % 墙内衰减
  \draw[thick,red!70!black] plot[domain=3:4.2,samples=30,smooth] (\x,{0.9+0.6*exp(-2.2*(\x-3))});
  \node[red!70!black] at (3.6,-0.5) {墙内指数衰减};
  % 透射波
  \draw[thick,green!50!black] plot[domain=4.2:8.2,samples=80,smooth] (\x,{0.9+0.18*sin(480*\x)});
  \node[green!50!black] at (6.4,1.9) {漏过去的残余振幅};
  \node at (4.2,-1.1) {隧穿：波函数在经典禁区内不骤降到零，薄墙外仍有非零概率};
\end{tikzpicture}
\end{center}

\begin{mathbridge}
完整的薛定谔方程（一维）写作
\[
  i\hbar\,\frac{\partial\psi}{\partial t}
  =\hat{H}\psi,\qquad
  \hat{H}=-\frac{\hbar^{2}}{2m}\frac{\partial^{2}}{\partial x^{2}}+V(x).
\]
左边是“状态随时间的变化率”；右边动能项里的二阶空间导数量度波函数的弯曲程度——弯得越急波长越短、动量越大，加上势能 \(V(x)\)，合起来就是“能量在空间和形状上的完整分布”。把自由粒子（\(V=0\)）的平面波 \(\psi=e^{i(kx-\omega t)}\) 代入，左边给出 \(\hbar\omega\)，右边给出 \(\hbar^{2}k^{2}/(2m)\)，正是 \(E=p^{2}/(2m)\)——方程在自由情形下精确归还了经典动能关系，约束三再次兑现。

找定态的办法是分离变量：令 \(\psi(x,t)=\varphi(x)\,e^{-iEt/\hbar}\)，代入后时间因子两边约去，剩下\emph{定态薛定谔方程}
\[
  \hat{H}\varphi=E\varphi .
\]
读法：寻找那些“被哈密顿量作用之后只变出一个倍数”的特殊波形，倍数 \(E\) 就是能量。盒中粒子代入边界 \(\varphi(0)=\varphi(L)=0\)，解得 \(\varphi_{n}(x)=\sqrt{2/L}\,\sin(n\pi x/L)\) 与 \(E_{n}=n^{2}h^{2}/(8mL^{2})\)——主线层的整数半波论证，在这里是严格的解。时间因子 \(e^{-iEt/\hbar}\) 的模恒为一，所以定态的 \(|\psi|^{2}\) 确实冻结：主线层的“整体空转”不是比喻，是因子分解的必然。概率守恒也可以从方程直接推出：由方程及其复共轭计算 \(\partial|\psi|^{2}/\partial t\)，结果总是一个“流量”的空间差——概率只在空间里搬来搬去，总量分毫不增不减，这正是左边那个 \(i\) 的功劳。
\end{mathbridge}

\begin{challenge}
为什么方程对时间是\emph{一阶}、还必须是\emph{复数}的？两条要求背后是同一个字：守恒。若对时间是二阶（像绳波方程），那么“此刻的波函数”不足以决定未来，还得额外交出“此刻波函数的变化率”——状态的概念就要膨胀，第~43~章“波函数即全部状态”的信条随之破产。一阶方程保证“知道现在就知道未来”，与牛顿力学“状态完备”的深刻传统一脉相承。而复数与 \(i\) 保证演化是\emph{幺正}的：形式解可以写成 \(\psi(t)=e^{-i\hat{H}t/\hbar}\psi(0)\)，演化算符 \(U(t)=e^{-i\hat{H}t/\hbar}\) 保持任何两个态之间的“重叠长度”不变，概率守恒是它的直接推论。这台机器还有一个值得记住的推论：一个物理量若与哈密顿量“对易”（第~46~章的语言），它的期望值就不随时间变化——能量、动量、角动量的守恒律在量子世界里换了新装，但一个都没少。最后交代一句边界预告：方程里动能写作 \(p^{2}/(2m)\)，这只是低速近似；当粒子速度接近光速，能量与动量的真实关系（第~40~章）插不进来，方程必须整体升级——那是第~49~章狄拉克方程的故事。
\end{challenge}

\step{模型的边界}

\begin{mainline}
\textbf{边界一：非相对论近似。}薛定谔方程把动能按 \(p^{2}/(2m)\) 记账，这在粒子速度远小于光速时极好，在接近光速时系统性出错。高能物理、重原子内层电子、反物质的存在，都超出了它的国境——由第~49~章的相对论性量子理论接管。

\textbf{边界二：哈密顿量与总能量不能无脑画等号。}本章反复说“能量是演化的发电机”，严格的表述是：\emph{哈密顿量}是演化的发电机。只有当势能不含时间、系统孤立时，哈密顿量才等于守恒的总能量。用随时间变化的外场驱动系统（比如用一束激光照射原子），哈密顿量照样生成演化，但“系统总能量守恒”不再成立——原子可以从光场吸收能量。把“哈密顿量”和“守恒的总能量”在任何场合混为一谈，是本方程最常见的概念性误用。

\textbf{边界三：波函数本身不可直接测量。}能进实验室账本的只有两类东西：模平方（概率分布）和相位\emph{差}（干涉条纹的位置）。整张波函数的整体相位永远测不到——所以“相位在旋转”是模型的语言，而“两个成分的相位差随时间拉开”才是可检验的事实。这条区分在理解定态时尤其要紧。

\textbf{边界四：测量时刻不归本章管。}薛定谔方程描述的是两次测量\emph{之间}的连续演化：光滑、可逆、概率守恒。测量那一刻的状态更新（所谓坍缩）是另一套规则，由第~45~章审理。把方程用到测量过程上、或指望它单独解释“为什么只出现一个结果”，都是越界。

再列四条本章特有的典型误用：

\begin{itemize}[leftmargin=1.8em]
  \item \emph{“定态就是粒子静止。”}错。冻结的是概率分布 \(|\psi|^{2}\)，不是运动。盒中基态的能量 \(E_{1}>0\) 在盒内（\(V=0\)）全部是动能——概率地形不动，地形下面暗流涌动。第三题选（b）。
  \item \emph{“隧穿是粒子在墙里挖了条路，墙内它的动能为负。”}错。前半句是名字的误导，已澄清；后半句犯了把经典记账硬塞进量子模型的错误——“墙内动能为负”这个量没有任何测量对应，你若真的进墙去测位置，就已经用一次相互作用改写了状态（第~45~章的规则），测到的也不再是原来的演化。
  \item \emph{“波包扩散就是粒子变胖。”}错。每次探测，电子仍然以整颗的面目出现在某个点上；变胖的是“可能找到它的区域”的分布，是一颗粒子可能性的铺开，不是一颗粒子身材的走样。
  \item \emph{“量子力学说能量必须离散。”}错。自由粒子的能量完全连续。离散永远来自边界条件——墙、轨道闭合、周期性。能级是“关押”的产物，不是量子的默认设置。
\end{itemize}
\end{mainline}

\step{回头解释开场现象}

\textbf{量子点电视。}对账已完成：把纳米颗粒当成盒子，\(E_{n}=n^{2}h^{2}/(8mL^{2})\) 给出平方反比的定价规则——颗粒从六纳米缩到三纳米，量子化的那部分能量抬到四倍，发光从红端推向蓝绿端。材料决定“底色”（固有能隙），尺寸决定“加价”。你客厅里那块屏幕，就是一行被点亮了一亿次的定态薛定谔方程。第一题（能量变四倍）和开场悬念，到此一并收回。

\textbf{永不变调的音符。}氢原子的红光为什么百亿年、百亿光年不走样？因为氢原子核外电子住在库仑势这个“天然的盒子”里（第~52~章会正式求解它），它的定态能级由同一个方程、同一个势唯一决定——全宇宙的氢原子解的是同一道题，答案当然一字不差。跃迁发出的光子能量等于两个能级之差，频率 \(f=\Delta E/h\)，于是光谱线成了宇宙中最稳定的刻度。锶原子钟把这种稳定性做成仪器：原子的定态不因旅途劳顿而“走音”，不像发条会疲劳、音叉会老化——定态的“不变”是方程层面的精确，不是工程层面的凑合。人类目前最精确的计时装置，抵押品正是本章那句话：\emph{能量确定的成分，相位只能匀速空转}。

\textbf{顺带收回三笔隧穿的账。}其一，你的 U 盘和固态硬盘：闪存把电子关在一层绝缘氧化物围成的“浮栅”里，那堵墙够厚，隧穿概率被指数压得极低——所以数据能保存十年量级；但也正因为是指数，墙一点点变薄、一点点受损，泄漏就会剧烈加快，这是闪存有擦写寿命、久置会丢数据的物理根源。其二，扫描隧道显微镜：让一根极细的针尖贴着样品表面扫过，针尖与样品间的真空就是一堵墙，隧道电流对间距指数敏感——间距每小零点一纳米，电流大约变化一个数量级——于是单个原子的高低起伏都纤毫毕现，人类第一次“看见”了原子排成的字。其三，太阳：太阳核心温度对应的质子热运动动能只有约一千电子伏，而两个质子靠近时的库仑排斥势垒高达百万电子伏量级——按经典判决，太阳根本点不燃；全靠隧穿那道指数曲线，质子才有微小但非零的概率“漏”到一起，聚变成氦。隧穿概率小，反而成就了太阳燃烧得又慢又稳——这盏亮了四十六亿年的灯，灯芯是一条指数衰减的波函数。

\step{脑洞实验与挑战题}

\begin{thought}[把人关进无限深势阱]
把一个 \(70\) 千克的人关进宽 \(1\) 米的刚性盒子，用 \(E_{1}=h^{2}/(8mL^{2})\) 算他的基态能量：
\[
  E_{1}=\frac{(6.63\times10^{-34})^{2}}{8\times70\times1^{2}}\approx8\times10^{-70}\ \text{J},
\]
对应的“不安分”速度 \(v=\sqrt{2E_{1}/m}\approx5\times10^{-36}\) 米每秒——以这个速度爬过一根头发丝的直径，需要约 \(10^{23}\) 年，是宇宙年龄的万亿倍。再看能级间隔：假如你随手翻一页书做了约 \(10^{-2}\) 焦的功，这笔钱相当于把这个人从基态一路激发到量子数 \(n\approx10^{34}\) 的能级——在那里，相邻能级的相对间隔只有 \(10^{-33}\) 量级，连续得不能再连续。这就是本节最极端的一课：量子化从未在宏观世界消失，它只是被 \(h^{2}\) 这个小到离谱的系数压进了任何仪器都够不着的缝隙里。经典力学不是错的，它是 \(E_{n}\) 密成一片之后看不清栅格的远观图。
\end{thought}

\begin{thought}[给太阳拧一拧隧穿的旋钮]
隧穿概率对势垒宽度指数敏感，而太阳核心质子恰好活在指数曲线的陡峭段上。开一个双方向脑洞。把透射概率调大一百万倍：聚变燃料将在远古时代就猛烈烧尽，恒星短寿，行星来不及冷却，更来不及长出看星星的生物。调小一百万倍：质子永远漏不过库仑墙，太阳根本不亮，地球冻成一块冰球。我们坐在一颗恰好亮了四十六亿年、还将再亮五十亿年的恒星旁边，部分原因正是那条指数曲线的斜率落在了“慢而稳”的区间里。同一个机制还在你自己的骨头边工作：烟雾报警器里的镅-241 靠 \(\alpha\) 粒子隧穿漏出原子核，漏出的速率恒定——放射性衰变的半衰期，正是“单位时间隧穿概率恒定”经过本章的时间演化积分出来的指数律。从下注到退休，隧穿这根指数曲线一直在悄悄管理世界的时间表。
\end{thought}

\begin{puzzles}
\item 盒中电子从 \(n=2\) 跃迁到 \(n=1\)，发出一个光子。现在把盒子放宽一倍，其他不变，发出的光波长变成原来的几倍？\emph{思路提示}：先写 \(\Delta E=E_{2}-E_{1}=3h^{2}/(8mL^{2})\)，看 \(L\) 出现在几次方上；再用光子能量与波长的关系换算。注意是波长，不是频率。
\item 把基态（\(E_{1}\)）和第一激发态（\(E_{2}\)）各占一半的盒中电子叠加起来，它出现在盒子左半边和右半边的概率会随时间此消彼长。这个“摇晃”的频率由什么决定？若 \(E_{2}-E_{1}\) 加倍，摇晃变快还是变慢？\emph{思路提示}：每个定态的相位以 \(\omega=E/\hbar\) 空转；概率里的交叉项只感受两个相位的\emph{差}。这和两个频率相近的音叉产生的拍是同一笔账。
\item 闪存工程师面对一个两难：绝缘层做厚，数据保持更久，但写入和擦除（同样靠隧穿注入和拉出电子）需要更高的电压、更慢的速度，还更容易损伤绝缘层。为什么“加厚一点点”会在两端都产生巨大的效果？\emph{思路提示}：关键词是“指数”。透射概率随厚度指数变化，意味着厚度改变百分之十，概率可能改变几个数量级——指数函数的世界里没有“温和的调整”。
\item 一位读者反驳：“盒中电子基态的概率分布不随时间变化，所以它动能为零；可公式 \(E_{1}=h^{2}/(8mL^{2})\) 又说它能量不为零，矛盾。”请指出他把哪两件事混为一谈了。\emph{思路提示}：盒内势能为零，\(E_{1}\) 是什么性质的能？再回想定态冻结的究竟是什么——是 \(|\psi|^{2}\)，还是波函数本身？一个模不变、相位飞转的复数，大小没变，它“没动”吗？
\end{puzzles}

至此，量子力学的主体机器已经装配完成：第~43~章给了状态（波函数）与读数规则（玻恩规则），本章给了状态在两次测量之间的发动机（薛定谔方程）。但机器运转时还有两个房间我们没有进去。一个是测量那一刻：连续演化如何与“只出现一个结果”衔接，状态更新是什么性质——那是第~45~章。另一个是这台机器的通用语言：一旦把波函数抽象成“向量”、把哈密顿量抽象成“算符”，定态方程 \(\hat{H}\varphi=E\varphi\) 会露出它的真面目——寻找本征值与本征向量，线性代数将成为量子力学的母语——那是第~46~章。而你本章见过的每一个例子：量子点、原子钟、闪存、太阳，都只是同一台机器换着场景空转或做功。
